%% kapa: exam=1 round=36 year=2025 period=1 form=A subject=civil_law subject_ko=민법 seq=21 number=21
%% source: Q-Net 감정평가사 기출문제 (cst00309) artlSeq=5231525 file=raw/1008/5231525/2025년 제36회 감정평가사 1차시험 1교시 문제지 원본.pdf
%% conversion: kordoc 4.24.51 -> tools/gen1.py (auto); review=proofread-sonnet; flags=none
\begin{question}{21}
물권에 관한 설명으로 옳지 않은 것은? (다툼이 있으면 판례에 따름)
\begin{choices}
  \item 물권의 객체는 물건에 한정되지 않는다.
  \item 법률 또는 관습법이 인정하지 않는 새로운 종류의 물권 창설은 허용되지 않는다.
  \item 미등기건물의 양수인은 소유권이전등기가 경료되지 않아도 사실상의 소유권을 취득한다.
  \item 근린공원을 자유롭게 이용할 수 있다는 사정만으로는 공원이용권이라는 배타적인 권리를 취득하였다고 할 수 없다.
  \item 분할의 절차를 거치지 않고 1필의 토지 중 일부에 관해 소유권보존등기를 경료할 수 없다.
\end{choices}
\end{question}
\answer{3}
